There would be many way to phrase it, possibly using topology. Instead we work internally to a hypothetical model satisfying duality for a generic non-commutative algebra where Zariski cover are surjections, and show a contradiction.

\begin{remark}
Assume $A$ a non-commutative algebra and $a:A$, then the ideal $(a)$ generated by $a$ consists of elements of the form $\Sigma_{i=1}^l b_iac_i$ where $b_i,c_i:A$.
\end{remark}

We give the key non-commutative algebra result.

\begin{lemma}\label{pushout-diagonal-matrices}
For all $i\in\{0,1\}$, we have a pushout square:
\[\begin{tikzcd}
R^2\ar[d,swap,"d"]\ar[r,"\pi_i"] & R\ar[d]\\
M_2(R)\ar[r] & 0
\end{tikzcd}\]
where the $\pi_i:R^2\to R$ is the $i$-th projections and $d$ is the diagonal map.
\end{lemma}

\begin{proof}
Consider an arbitrary commutative diagram:
\[\begin{tikzcd}
R^2\ar[d,swap,"d"]\ar[r,"\pi_i"] & R\ar[d,"g"]\\
M_2(R)\ar[r,swap,"f"] & A
\end{tikzcd}\]
and let us prove $A=0$. Consider the unique $j\in\{0,1\}$ such that $j\not=i$. Then $d(e_j) = e_{j,j}$, and since $1 = e_{0,0}+e_{1,1} = \sum_{k=0}^1 e_{k,j}e_{j,j}e_{j,k}$ in $M_2(R)$, we have that $(e_{j,j}) = M_2(R)$. Then $(f(e_{j,j})) = A$, but we have that $f(e_{j,j}) =f(d(e_j)) = g(\pi_i(e_j)) = g(0) = 0$. So $0=_A1$ and $A=0$.
\end{proof}

\begin{proposition}
Assume a non-commutative ring $R$ with duality for f.p.\ non-commutative $R$-algebras. If the canonical map $1+1\to \Spec(R^2)$ is surjective, then $0=_R1$.
\end{proposition}

\begin{proof}
By \Cref{pushout-diagonal-matrices} and duality, we have a pullback square:
\[\begin{tikzcd}
(0=_R1) + (0=_R1)\ar[d]\ar[r] & 1+1\ar[d]\\
\Spec(M_2(R))\ar[r] & \Spec(R^2)
\end{tikzcd}\]
The right map is assumed surjective, so the left map is surjective. So $\Spec(M_2(R))\to\propTrunc{(0=_R1) + (0=_R1)}$, giving $\Spec(M_2(R))\to \Spec(0)$. Through duality we get $M_2(R) = 0$, therefore $0=_R1$.
\end{proof}